\documentclass[12pt,reqno]{amsart}
\usepackage[a4paper, hmargin={2.7cm,2.7cm},vmargin={3.3cm,3.3cm}]{geometry}
\usepackage[T1]{fontenc}
\usepackage[utf8]{inputenc}
\usepackage{lmodern,amsmath,amssymb,amsthm,mathtools,microtype}
\usepackage[hidelinks]{hyperref}
\mathtoolsset{showonlyrefs}
\newtheorem{theorem}{Theorem}[section]
\newtheorem{lemma}[theorem]{Lemma}
\newtheorem{proposition}[theorem]{Proposition}
\newtheorem{corollary}[theorem]{Corollary}
\theoremstyle{definition}
\newtheorem{definition}[theorem]{Definition}
\newtheorem{remark}[theorem]{Remark}
\newcommand{\R}{\mathbb R}
\newcommand{\Sph}{\mathbb S^2}
\newcommand{\Sym}{\operatorname{Sym}}
\newcommand{\tr}{\operatorname{tr}}
\newcommand{\rank}{\operatorname{rank}}
\newcommand{\im}{\operatorname{im}}
\newcommand{\norm}[1]{\lVert#1\rVert}
\newcommand{\ip}[2]{\langle#1,#2\rangle}
\newcommand{\Eh}{\widehat E}
\newcommand{\dd}{\mathrm d}
\numberwithin{equation}{section}
\setlength{\emergencystretch}{2em}
\raggedbottom
\title[Optimal eight-point configurations]{Optimal eight-point configurations for all nonnegative exponents}
\author{}
\date{September 18, 2026}
\hypersetup{pdftitle={Optimal eight-point configurations for all nonnegative exponents}}
\begin{document}
\begin{abstract}
We prove, with computer assistance, that the height-optimized square
antiprism uniquely minimizes logarithmic energy and Riesz energy of every
positive exponent among eight-point configurations on the unit sphere,
up to congruence.
\end{abstract}
\maketitle

\section{Introduction and results}
For eight distinct points $Y=(y_1,\ldots,y_8)$ on the unit sphere
$\Sph\subset\R^3$, define
\begin{equation}\label{eq:energy}
 E_s(Y)=\sum_{i<j}|y_i-y_j|^{-s}\quad(s>0),\qquad
 E_0(Y)=-\sum_{i<j}\log|y_i-y_j|.
\end{equation}
At a collision the energy is $+\infty$. Congruence means equality after
an orthogonal transformation and a permutation of the points. The
notation $s=0$ always refers to logarithmic energy. The direct
substitution $s=0$ in the Riesz formula would instead give a constant.

For $0<a<1$, let $X(a)$ be the square antiprism
\begin{equation}\label{eq:family}
 X(a)=\left\{\left(\sqrt{1-a}\cos\frac{j\pi}{4},
 \sqrt{1-a}\sin\frac{j\pi}{4},(-1)^j\sqrt a\right):0\le j<8\right\}.
\end{equation}
The parameter $a$ is the squared height. For every $s\ge0$ there is a
unique minimizer $a_s$ of $a\mapsto E_s(X(a))$; an exact equation for it
is given below.

\begin{theorem}\label{thm:main}
Let $s\ge0$. For every eight-point configuration
$Y$ on $\Sph$,
\[
 E_s(Y)\ge E_s(X(a_s)).
\]
Equality holds if and only if $Y$ is congruent to $X(a_s)$.
\end{theorem}

The square antiprism changes with $s$. Thus the theorem concerns a
one-parameter family of optimal configurations. The logarithmic member
has $a_0=(2\sqrt{58}-13)/7$.

For bounded exponents we use the three-point method of Cohn and
Woo~\cite{CW}. We construct a polynomial minorant of the pair
potential and an exact three-point identity whose summed terms are
nonnegative. Strict positivity of the certificate matrices also
determines the equality case. For large exponents we first locate every
possible global minimizer near the packing antiprism. We then place it
and the optimized energy antiprism in a common convex coordinate set
on which the full energy has positive Hessian.

The integer arithmetic checks are described in
Section~\ref{sec:arithmetic}. The accompanying programs and certificates
verify its finite computational obligations; their use is described in
\path{CAP.md}. The exact exponent intervals and verification records are
included with the code.

\subsection{The optimized height}
It is convenient to use the normalized energy
\[
 \Eh_s(Y)=\sum_{i<j}\phi_s(y_i\cdot y_j),\qquad
 \phi_s(t)=
 \begin{cases}((2-2t)^{-s/2}-1)/s,&s>0,\\
 -\tfrac12\log(2-2t),&s=0.
 \end{cases}
\]
Thus $\Eh_s=(E_s-28)/s$ for $s>0$, and the minimizers are unchanged.
The four squared distances of $X(a)$ are
\begin{equation}\label{eq:distances}
\begin{aligned}
 d_A&=2(1-a),& d_B&=4(1-a),\\
 d_C&=2-\sqrt2+(2+\sqrt2)a,&
 d_D&=2+\sqrt2+(2-\sqrt2)a,
\end{aligned}
\end{equation}
with multiplicities $8,4,8,8$, respectively. Write
$b_-=2+\sqrt2$, $b_+=2-\sqrt2$.

\begin{lemma}[Exact optimized parameter]\label{lem:parameter}
For every $s\ge0$ there is exactly one $a_s\in(0,1)$ satisfying
\begin{equation}\label{eq:root}
 F(s,a_s)=0,
\end{equation}
where, with $q_s=1+s/2$,
\begin{equation}\label{eq:F}
 F(s,a)=16d_A^{-q_s}+16d_B^{-q_s}
        -8b_-d_C^{-q_s}-8b_+d_D^{-q_s}.
\end{equation}
For $s>0$, $a_s$ uniquely minimizes $G_s(a)=E_s(X(a))$. At zero,
\begin{equation}\label{eq:a0}
 a_0=\frac{2\sqrt{58}-13}{7}.
\end{equation}
The normalized family minimum and the parameter depend smoothly on $s$
near zero and throughout the positive half-line.
\end{lemma}
\begin{proof}
Differentiating the normalized family energy gives
$\partial_a\Eh_s(X(a))=F(s,a)/2$, including at $s=0$ by smooth extension.
Furthermore,
\begin{equation}\label{eq:Fa}
 \partial_aF=q_s\left(32d_A^{-q_s-1}+64d_B^{-q_s-1}
             +8b_-^2d_C^{-q_s-1}+8b_+^2d_D^{-q_s-1}\right)>0.
\end{equation}
At $a=0$, the two positive terms of $F$ sum to at most $12$, whereas
the magnitude of its negative $d_C$ term is greater than
$8(2+\sqrt2)>12$. As $a\uparrow1$, the positive terms tend to infinity
and the negative terms remain bounded. Existence and uniqueness follow.
At $s=0$, direct simplification gives
\[
 \frac12F(0,a)=\frac{2(7a^2+26a-9)}{(1-a)(a^2+6a+1)},
\]
which gives \eqref{eq:a0}. The implicit-function theorem applies because
\eqref{eq:Fa} is strictly positive. The smooth extension of the normalized
potential follows, for example, from
\[
 \phi_s(t)=\ell(t)\int_0^1 e^{s\ell(t)u}\,du,
 \qquad \ell(t)=-\tfrac12\log(2-2t).
\]
\end{proof}





The family minimizer is stationary for the energy on $(\Sph)^8$.
Indeed, the symmetries of the antiprism act transitively on its vertices.
A reflection fixing a vertex reverses its azimuthal tangent direction,
so the energy gradient there has only a meridional component.
Transitivity identifies these components with the derivative of the
height variation. That derivative vanishes at $a_s$, so every tangential
gradient vanishes.

\section{A certificate for a bounded exponent interval}\label{sec:finite}
We give one criterion for all the finite certificates. Let $I\subset[0,\infty)$
be a closed interval and $d\ge6$ an integer. All definitions in this
section are made for each $s\in I$ at the exact root $a=a_s$.
Put $e_s=\Eh_s(X(a_s))$ and
\begin{equation}\label{eq:nodes}
 A=a,\quad B=2a-1,\quad
 C=-a+(1-a)/\sqrt2,\quad D=-a-(1-a)/\sqrt2.
\end{equation}
We require $0<a<1/3$ and that these nodes are distinct, nonzero, in
$(-1,1)$, and have no opposite pair. The unordered triangle types of
$X(a)$ are
\[
 (D,D,A),\quad(D,B,C),\quad(D,C,A),\quad(B,A,A),\quad(C,C,A),
\]
with respective multiplicities $8,16,16,8,8$. Their permutations give 21 distinct
ordered triples. These statements follow directly from \eqref{eq:family}.

\subsection{A scalar minorant}
Choose positive even contact multiplicities $m_A,m_B,m_C,m_D$, and a
positive odd endpoint multiplicity $m_-$, with
$m_-+m_A+m_B+m_C+m_D=2d+1$.
Let $H_s$ be the polynomial of degree at most $2d$ whose derivatives
through order $m_\alpha-1$ agree with those of $\phi_s$ at each node
$\alpha\in\{-1,A,B,C,D\}$. Counting zeros with multiplicity proves
uniqueness, and hence existence by finite-dimensional linear algebra.

\begin{lemma}\label{lem:minorant}
For $-1\le t<1$, $H_s(t)\le\phi_s(t)$, with equality precisely at
$-1,A,B,C,D$.
\end{lemma}
\begin{proof}
For $j\ge1$,
\[
 \phi_s^{(j)}(t)=2^{j-1}\prod_{k=1}^{j-1}(s/2+k)
                    (2-2t)^{-s/2-j}>0,
\]
including $s=0$. Away from the interpolation nodes, the confluent
interpolation remainder is a positive derivative divided by $(2d+1)!$,
multiplied by
$(t+1)^{m_-}\prod_{\alpha=A,B,C,D}(t-\alpha)^{m_\alpha}$.
This product is strictly positive for $t>-1$ away from the contacts.
At the nodes equality is part of the interpolation conditions.
\end{proof}
Consequently, writing $E_H(Y)=\sum_{i<j}H_s(y_i\cdot y_j)$,
\begin{equation}\label{finite73:scalarenergy}
 \Eh_s(Y)\ge E_H(Y),\qquad E_H(X(a_s))=e_s.
\end{equation}
We allow a positive scalar multiplier $\chi(s)$ of the certificate
identity. The archived cells use either $\chi=1$ or, for a positive
cell center $c$, $\chi(s)=(s/c)(3/2)^{(s-c)/2}$ on a positive interval.

\subsection{Moment and contact-kernel polynomials}
For a triple of points use coordinates $(u,v,t)=(x\cdot y,x\cdot z,y\cdot z)$. Let $\operatorname{Sym}$ mean the average over all six permutations of these coordinates, including repeated terms. Set
\begin{equation}\label{finite73:Q}
 Q_0=1,\quad Q_1=t-uv,\quad
 Q_k=2(t-uv)Q_{k-1}-(1-u^2)(1-v^2)Q_{k-2}.
\end{equation}
For $0\leq k\leq d-1$ and $0\leq i,j\leq d-k$, form the symmetric matrix polynomial
$S_{k,ij}=\operatorname{Sym}(u^iv^jQ_k)$ and put
\begin{equation}\label{finite73:R}
 R_k=6S_k+S_k(u,u,1)+S_k(v,v,1)+S_k(t,t,1)
       +\frac{\delta_{k0}}7\mathbf J,
\end{equation}
where $\mathbf J$ is the all-ones matrix of the required order. Let $M_k(Y)$ be the sum of $S_k$ over all ordered triples of labels, including repeated labels, and let $\Psi_Y$ sum over ordered triples with distinct labels.

\begin{lemma}[Moment positivity and counting]\label{finite73:moment}
For every spherical configuration, $M_k(Y)\succeq0$ and $\Psi_Y(R_k)=6M_k(Y)$.
\end{lemma}
\begin{proof}
At a fixed first vertex $x$, the recurrence is the polynomial version of
$((1-u^2)(1-v^2))^{k/2}\cos(k\theta)$ for the angle between the projections of $y,z$ into $x^\perp$. Its double sum is a sum of cosine and sine feature outer products. For $k=0$ it is the usual moment-vector outer product. If a tangent projection vanishes and $k>0$, the prefactor is zero; hence the polynomial formula covers these cases continuously. Summing in $x$ gives positivity, and averaging permutations does not change the all-ordered sum.

Separate all ordered triples into those with distinct labels, exactly two equal labels, and three equal labels. Every unordered pair belongs to six unordered triples. The three repeated-coordinate terms in \eqref{finite73:R}, when summed over ordered distinct triples, therefore supply six times the exactly-two-equal contribution to $M_k$. The all-equal contribution is $8\delta_{k0}\mathbf J$, while $\Psi_Y(\delta_{k0}\mathbf J/7)=336\delta_{k0}\mathbf J/7=48\delta_{k0}\mathbf J$. This proves the identity.
\end{proof}

For the candidate, reflection fixing a vertex makes each sine moment vanish, while vertex transitivity identifies the remaining cosine moment vectors up to sign. Thus each $M_k(X(a_s))$ has rank at most one. A verified positive diagonal pivot makes its rank exactly one. If its pivot index is $r$, its kernel is generated by
\begin{equation}\label{finite73:momentkernel}
 e_f-\frac{M_k[r,f]}{M_k[r,r]}e_r\qquad(f\ne r).
\end{equation}
Multiplication of this generator matrix by the fixed witness matrix $J_k$ gives $N_k$. These columns still belong to the exact kernel. Only the $k=0$ matrix $J_0$ needs a separate invertibility check for the equality proof below.

For the three SOS blocks, take the raw monomial columns of degrees at most $d,d-1,d-2$ respectively. Evaluate at the $21$ ordered contacts. A verified invertible $21$-column pivot subsystem defines exact kernel generators with identity free rows and pivot rows $-E_{\rm piv}^{-1}E_{\rm free}$. Multiply these generators by their fixed witness change-of-coordinates matrices. Denote the resulting polynomial feature columns by $z_0,z_1,z_\Delta$. Each feature vanishes at all contacts. Without a further invertibility test on its fixed change of coordinates we assert kernel membership, not that every transformed list is a full basis.

\subsection{The exact coefficient identity}
Set
\begin{equation}\label{finite73:target}
\begin{aligned}
 P_s(u,v,t)&=\frac{H_s(u)+H_s(v)+H_s(t)}3-\frac{e_s}{28},\\
 \Delta(u,v,t)&=1+2uvt-u^2-v^2-t^2.
\end{aligned}
\end{equation}
The required identity, as an identity of polynomial coefficients, is
\begin{align}\label{finite73:identity}
 \chi(s)P_s={}&\sum_{k=0}^{d-1}\operatorname{tr}(F_kR_k)
 +\operatorname{Sym}(z_0^TW_0z_0)
 +\operatorname{Sym}((1-u^2)z_1^TW_1z_1)\notag\\
 &+\Delta\operatorname{Sym}(z_\Delta^TW_\Delta z_\Delta),
 \qquad F_k=N_kB_kN_k^T,
\end{align}
where all reduced symmetric matrices $B_k,W_0,W_1,W_\Delta$ are positive definite. Every term has degree at most $2d$: the determinant-localized term even has degree at most $2d-1$. On a realized sphere triple, $1-u^2\geq0$ and $\Delta\geq0$ since $\Delta$ is the Gram determinant.

Each unordered pair occurs twelve times in the ordered-distinct sum of $(H(u)+H(v)+H(t))/3$, and there are $336$ ordered distinct triples. Thus
\begin{equation}\label{finite73:configurationidentity}
 \Psi_Y(\chi P_s)=12\chi(s)(E_H(Y)-e_s)\geq0.
\end{equation}
Indeed the SOS terms are nonnegative individually, and the summed moment term is $6\operatorname{tr}(F_kM_k(Y))\geq0$ by Lemma~\ref{finite73:moment}. We do not require a pointwise positivity assertion for $R_k$. Equations~\eqref{finite73:scalarenergy} and \eqref{finite73:configurationidentity} imply the global energy inequality for every configuration; competitors need not be stationary or symmetric.

Here is why a corrected identity on the retained rows is an exact identity on all rows. Let $V_{2d}$ be the space of symmetric three-variable polynomials of degree at most $2d$. The coefficient coordinates are the coefficients of individual canonical sorted monomials, not sums over their permutation orbits. Put $m=\dim V_{2d}$ and define
\begin{equation}\label{finite73:annihilators}
 L_0(f)=\Psi_{X(a_s)}f,\qquad
 L_1(f)=\left.\frac{d}{da}\Psi_{X(a)}f\right|_{a=a_s}.
\end{equation}
In the second expression all coefficients of $f$, and in particular all candidate kernel matrices, $s$, and $\chi(s)$, are frozen. The pair sum of squared inner products in $X(a)$ is $48a^2-32a+12$. Consequently
\begin{equation}\label{finite73:annihilatorindependence}
 L_0(1)=336,\quad L_1(1)=0,\quad
 L_1(u^2+v^2+t^2)=1152(3a_s-1)\ne0.
\end{equation}
The common kernel $K=\ker L_0\cap\ker L_1$ therefore has dimension $m-2$.

Every coefficient column on the right side of \eqref{finite73:identity} belongs to $K$. A product of two contact-vanishing features has both value and first height derivative zero, even after multiplication by a localizer. For a moment column, freeze $N_k$ at $a_s$. The matrix $N_k^TM_k(X(a))N_k$ is positive semidefinite for nearby $a$ and is zero at $a_s$. Every quadratic form has a differentiable minimum there, so its derivative vanishes; polarization gives zero for every entry. This proves the assertion using the counting identity. The target also belongs to $K$: its value is zero by interpolation and $e_s=\Eh_s(X(a_s))$, and its derivative is $12\chi(s)\partial_a\widehat E_s(X(a_s))=6\chi(s)F(s,a_s)=0$, since Hermite first derivatives agree with $\phi_s'$.

Collect the independent symmetric matrix entries in $w$, with the explicit convention
\begin{equation}\label{finite73:svec}
 w_{ii}=B_{ii},\qquad w_{ij}=\sqrt2 B_{ij}\quad(i<j).
\end{equation}
Unpacking divides off-diagonal coordinates by $\sqrt2$. Let $A(s)w=b(s)$ be the full coefficient system for \eqref{finite73:identity}. Its columns and target lie in $K$. The witness selects $m-2$ rows $R$, an $n\times(m-2)$ correction matrix $C_r$, an initial vector $w_0$, and optionally a square left preconditioner $L$ (otherwise $L=I$). Define
\begin{equation}\label{finite73:correction}
 T=L A_R C_r,\quad r=L(b_R-A_Rw_0),\qquad
 w=w_0+C_rT^{-1}r.
\end{equation}
If $T$ is invertible, $L$ is invertible and $A_R$ has rank $m-2$. Since $\operatorname{im}A\subseteq K$, it follows that $\operatorname{im}A=K$, and restriction to the selected rows is injective on $K$. Equation~\eqref{finite73:correction} makes the selected residual zero. The full residual lies in $K$, so every coefficient is zero. No assumption about a particular distinguished pair of omitted rows is required.



\begin{proposition}[Finite certificate criterion]\label{finite73:obligation}
Suppose the node conditions above hold throughout $I$. Suppose that the
moment pivots are positive, the contact evaluation pivots and the
coefficient correction matrix $T$ are invertible, and every corrected
matrix in \eqref{finite73:identity} is positive definite throughout $I$.
Suppose also that $J_0$ is invertible and $z_0$ is nonzero at each of
$(A,A,A),(A,A,C),(A,A,D)$. Then $X(a_s)$ uniquely minimizes $E_s$ up to
congruence for every $s\in I$, with the logarithmic convention at zero.
\end{proposition}
\begin{proof}
The exact coefficient correction proves \eqref{finite73:identity}.
The moment counting identity and the nonnegative localizers give
\eqref{finite73:configurationidentity}. Combining it with
\eqref{finite73:scalarenergy} proves the energy inequality.
The following argument proves uniqueness.
\end{proof}

\subsection{Equality classification}
\begin{lemma}[Equality from the finite certificates]\label{finite73:equality}
Under the hypotheses of Proposition~\ref{finite73:obligation}, if $E_s(Y)=E_s(X(a_s))$, then $Y$ is $X(a_s)$ up to an orthogonal transformation and relabeling.
\end{lemma}
\begin{proof}
Equality in the scalar comparison and in \eqref{finite73:configurationidentity} forces every off-diagonal inner product to lie in $\{-1,A,B,C,D\}$. There is no antipodal pair: for such a pair and any third distinct vertex its two inner products would be $t$ and $-t$ in $\{A,B,C,D\}$. The nonzero and no-opposite-pair checks exclude this, including the possibility $t=0$.

For each vertex $y$, put $m_y=(\sum_z(y\cdot z)^j)_{j=0}^{d}$, where the sum includes the vertex itself. The $k=0$ full moment is $M_0(Y)=\sum_y m_y m_y^T$. The candidate has $M_0(X(a_s))=8m_*m_*^T$ for its common moment vector $m_*$. Since the $B_0$ block is positive definite and its nonnegative summed term vanishes, $N_0^Tm_y=0$ for every $y$. The verified invertibility of $J_0$ makes the columns of $N_0$ span $m_*^\perp$. Hence $m_y$ is a scalar multiple of $m_*$. Their zeroth coordinates are both $8$, so $m_y=m_*$.

Subtracting the self contribution at inner product $1$, the moments of orders $0,1,2,3$ form a Vandermonde system in the four distinct contact nodes. They force the per-vertex neighbor counts
\begin{equation}\label{finite73:neighbors}
 n_A=2,\qquad n_B=1,\qquad n_C=2,\qquad n_D=2.
\end{equation}
The first moment of the candidate is zero, so $y\cdot\sum_z z=0$ for every $y$. Summing in $y$ gives $|\sum_z z|^2=0$, proving that the configuration has zero centroid.

The unweighted $W_0$ term in \eqref{finite73:identity} is strictly positive definite on nonzero feature vectors. Since the total sum vanishes, $z_0$ must vanish at each equality triangle. The verified nonzero feature coordinates at $(A,A,A),(A,A,C),(A,A,D)$ forbid those triangle types. Thus two incident $A$ edges force their opposite edge to have inner product $B$.

The simple graph of $A$ edges is two-regular. A triangle is forbidden. A cycle of length at least five would give every vertex two distinct distance-two neighbors, both $B$ neighbors, contradicting \eqref{finite73:neighbors}. Therefore the graph is two four-cycles. For each, the Gram matrix in cyclic order is
\begin{equation}\label{finite73:squaregram}
 \begin{pmatrix}
 1&a&2a-1&a\\a&1&a&2a-1\\2a-1&a&1&a\\a&2a-1&a&1
 \end{pmatrix}.
\end{equation}
Its row sums are $4a$. If $c_*$ is this four-cycle's centroid, then $|c_*|^2=a$ and each vertex has inner product $a$ with $c_*$. Subtracting $c_*$ gives four vectors of squared norm $1-a$, adjacent inner products zero, and opposite vectors equal to negatives. They are a regular square in $c_*^\perp$. The total centroid condition makes the other square's centroid $-c_*$, so the squares have a common axis. All cross inner products are $C$ or $D$, since the $A$ and $B$ neighbors are exhausted within each square. Their planar angle cosines are consequently $\pm1/\sqrt2$, which forces an angular offset $\pi/4$ modulo $\pi/2$. This is precisely \eqref{eq:family}.
\end{proof}



\subsection{Application to the finite range}
The verified closed cells cover $[0,124]$. Each satisfies the hypotheses
of Proposition~\ref{finite73:obligation}; the root tubes, inverse
residuals, strict matrix bounds and equality exclusions are retained in
the machine-readable reports. The exact endpoints and the witness and
source bindings are the replay inputs, not rounded numbers read from a
table. The independent aggregation recomputes their union using rational
arithmetic and rejects a missing or failed cell. This proves
Theorem~\ref{thm:main} for $0\le s\le124$.

The scalar interpolation pattern is allowed to change between cells.
No agreement between the two auxiliary polynomials is needed at a shared
endpoint: either complete certificate suffices there.

\section{The packing configuration}\label{r3:packing}
Put
\[
 a=\frac{2\sqrt2-1}{7},\qquad P=X(a),\qquad \tau=2(1-a),\qquad
 T(Y)=\max_{i<j}y_i\cdot y_j.
\]
A contact is a pair attaining $T$. We call a configuration \emph{contact-minimal
locally optimal} if it has a neighborhood containing neither a smaller value
of $T$ nor a configuration with the same $T$ and fewer contacts.

\begin{theorem}\label{r3:packingclassification}
Every contact-minimal locally optimal eight-point configuration with
$T<3/10$ is congruent to $P$. Moreover, $T(Y)\ge a$ for all configurations,
with equality exactly on the congruence orbit of $P$.
\end{theorem}

\subsection{Contact graph and facial geometry}
Write $t=T(Y)<3/10$. The Gram matrix $G$ has trace $8$, is positive
semidefinite, and has rank at most three. For $-1\le u\le t$,
$u^2\le(t-1)u+t$. Hence, with $S=\sum_i y_i$,
\[
 \frac{64}{3}\le\tr G^2\le16+48t+(t-1)\norm S^2\le16+48t,
\]
so $t\ge1/9$. Contact edges have spherical length $\theta=\arccos t$,
with $2\pi/5<\theta<\pi/2$. Their shorter great-circle arcs cannot cross:
from each of two crossing edges choose an endpoint within $\theta/2$ of
the crossing. The strict spherical triangle inequality would give a pair
at distance less than $\theta$. Overlapping arcs on one great circle
likewise give a too-close pair of endpoints.

Two neighbors of a vertex have tangent angle at least
\[
 \alpha=\arccos\frac{t}{1+t},\qquad 75^\circ<\alpha<90^\circ.
\]
Thus all degrees are at most four. If all contact tangent directions at
one vertex lie in a closed semicircle, choose a unit tangent vector $w$
with nonpositive inner products with them. Along the motion
$x(u)=\cos u\,x+\sin u\,w$, every contact inner product becomes
$t\cos u+(w\cdot y)\sin u<t$ for sufficiently small $u>0$.
Every strict noncontact retains its slack. This contradicts the local
selection condition. In particular degrees one and two are impossible,
and every facial angle at an active vertex is less than $\pi$.

A nontrivial component has at least five vertices: a four-vertex component
of minimum degree three would be $K_4$, whose contact Gram matrix
$(1-t)I+tJ$ has rank four. There is therefore exactly one active component.
It is two-connected. Indeed, if a cut vertex separated off at most three
vertices, minimum degree three would be impossible or would force a
contact $K_4$ including the cut vertex. Every component after deletion of
a cut vertex would thus have at least four vertices, requiring at least
nine vertices in total. Consequently the active component has simple
facial cycles bounding disks on the sphere.

We supply the global perimeter fact needed to restrict these faces.
\begin{lemma}[One-turn polygon]\label{r3:topology-one-turn}
Let $p_0,p_1,\ldots,p_{n-1},p_n=p_0$ be a closed polygonal
chain in an oriented Euclidean plane, with all edges nonzero. Suppose
the edge directions have arguments $\vartheta_0<\cdots<\vartheta_{n-1}
<\vartheta_0+2\pi$ and every successive difference, including
$\vartheta_0+2\pi-\vartheta_{n-1}$, lies in $(0,\pi)$.
Then the chain is the positively oriented boundary of a strictly convex
polygon. Equivalently, every vertex not on a given edge lies strictly
to that edge's left.
\end{lemma}
\begin{proof}
Fix an edge, rotate and cyclically relabel so that it is the first edge
and has argument zero, and write $H_j$ for the signed height of $p_j$
above its line. Thus $H_0=H_1=H_n=0$. For $j=1,\ldots,n-1$,
\[
 H_{j+1}-H_j=|p_{j+1}-p_j|\sin\vartheta_j.
\]
The increments are first positive, possibly then zero for one edge,
and then negative. The first is strictly positive because
$0<\vartheta_1<\pi$; the last is strictly negative because
$\pi<\vartheta_{n-1}<2\pi$. Consequently $H_j>0$ for
$2\le j\le n-1$: before the sign change this follows from the first
increment, and after it from the negative increments remaining before
$H_n=0$. Each edge therefore exposes precisely its two endpoints in
the convex hull of all vertices. The vertices are distinct, all the
displayed segments are successive exposed edges, and the chain is the
boundary of that convex hull, oriented with its interior on the left.
\end{proof}

\begin{lemma}[Spherical perimeter]
\label{r3:perimeter}
A simple spherical polygon bounding a disk on its left, with short
geodesic edges and all interior angles in $(0,\pi)$, has perimeter
less than $2\pi$.
\end{lemma}
\begin{proof}
Let its vertices be $x_i$, its side lengths $\theta_i\in(0,\pi)$,
and its interior angles $\alpha_i\in(0,\pi)$. Set
$n_i=(x_i\times x_{i+1})/\sin\theta_i$. The short arc from
$n_{i-1}$ to $n_i$ has length $\pi-\alpha_i$ and left normal $x_i$.
By Gauss--Bonnet on the original Jordan disk $D$, this closed polar
chain has length
\[
 L_*=\sum_i(\pi-\alpha_i)=2\pi-\operatorname{area}(D)<2\pi.
\]
No simplicity of the polar chain is being assumed.

Every closed spherical curve of length $L<2\pi$ lies in an open
hemisphere. Indeed, choose points $U,V$ separated by half its
arclength. For any point $z$ on either half, spherical distances
$p=d(U,z)$ and $q=d(V,z)$ satisfy $p+q\le L/2<\pi$, and hence
\[
 (U+V)\cdot z
 =2\cos((p+q)/2)\cos((p-q)/2)>0.
\]
In particular $U+V\ne0$, and its direction is a hemisphere center.
Apply this observation to the polar chain and denote its hemisphere
center by $c$.

The original boundary avoids both $c$ and $-c$, since its edge planes
have normals $n_i$ with $c\cdot n_i>0$. Along each edge its longitude
$\varphi$ about $c$ has derivative
\[
 \varphi'=\frac{c\cdot n_i}{1-(c\cdot x)^2}>0
\]
when the edge is parametrized by spherical arclength. Its total
longitude change is therefore $2\pi d$ for an integer $d\ge1$.
Choose a meridian from $c$ to $-c$ avoiding all vertices. It has
exactly $d$ transverse intersections with the boundary, all with the
same crossing orientation. On a Jordan disk, successive transverse
crossings of a simple path alternate between entrance and exit, and
thus have opposite orientations. It follows that $d\le1$, hence
$d=1$.

Gnomonically project the polar chain to $c^\perp$ by
$n\mapsto n/(c\cdot n)-c$. Its edge corresponding to $x_i$ has
nonzero left normal proportional to
$X_i=x_i-(c\cdot x_i)c$. The arguments of these normals increase
exactly once around the circle by the preceding longitude argument.
Moreover
\[
 \det_c(X_i,X_{i+1})
 =c\cdot(x_i\times x_{i+1})>0.
\]
The longitude change on this edge is in $(0,\pi)$: its short arc is
a normalized positive combination of its endpoints, whose projections
lie in that open planar sector. Thus every successive normal turn
is in $(0,\pi)$. Lemma~\ref{r3:topology-one-turn}, applied equivalently
to edge directions, shows that the projected polar chain is a simple
strictly convex polygon. Its inverse image bounds a positive-area
spherical disk $D^*$ in the hemisphere centered at $c$.

At polar vertex $n_i$, the incoming and outgoing left normals are
$x_i$ and $x_{i+1}$, both tangent to the sphere at $n_i$. Their
positive angle is exactly $\theta_i$, since
$n_i=(x_i\times x_{i+1})/\sin\theta_i$. These are therefore the
exterior angles of $D^*$. A second application of Gauss--Bonnet gives
\[
 \sum_i\theta_i=2\pi-\operatorname{area}(D^*)<2\pi.
\]
\end{proof}

Every active face therefore has three or four edges, since
$\theta>2\pi/5$. A contact triangle has a small convex region containing
no additional point: a nonvertex in it is a normalized positive combination
$v=(\sum c_i x_i)/N$, with $C=\sum c_i>N$; because $t>0$,
$v\cdot x_j=[tC+(1-t)c_j]/N>t$. Its complementary region contains other
active vertices, so a triangular face uses the small region.

For a contact four-cycle $x_1,x_2,x_3,x_4$, put
$r=x_1\cdot x_3$, $s=x_2\cdot x_4$. The differences
$x_1-x_3,x_2-x_4$ are nonzero and orthogonal, and both are perpendicular
to the two sums $x_1+x_3,x_2+x_4$. The sums are therefore parallel; they
point in the same direction because their inner product is $4t>0$.
In particular
\[
 (1+r)(1+s)=4t^2.
\]
Writing their common unit direction as $m$, the four corners have heights
$h_1=\sqrt{(1+r)/2}$ and $h_2=\sqrt{(1+s)/2}$, with $h_1h_2=t$.
Because $r,s\le t$, both heights are at least
$t\sqrt{2/(1+t)}>t$. Gnomonic projection about $m$ gives a convex diamond.
For a convex combination $w$ of the corners, $\norm w\ge m\cdot w>t$.
Its normalized direction cannot have inner product at most $t$ with all
corners, since that would imply $\norm w\le t$ by averaging.
The small quadrilateral is therefore empty of other packing points,
and every quadrilateral face uses that region. An isolated point can
lie neither in any face, nor on an edge, nor at a vertex. All eight
vertices are consequently active.

Let $A,B$ be the alternating angles of a contact quadrilateral. Opposite
angles agree by the spherical cosine law, and
\begin{equation}\label{pack:rhombus}
 \cos A=\frac{s-t^2}{1-t^2},\quad
 \cos B=\frac{r-t^2}{1-t^2},\quad
 \cot(A/2)\cot(B/2)=t.
\end{equation}
The last equality follows by squaring, using the diagonal identity, and
then using positivity of the half-angle cotangents. For a face of the
\emph{exact} contact graph the diagonals are strict noncontacts. Thus
\begin{equation}\label{pack:angles}
 \alpha<A,B<2\alpha,
\end{equation}
since the decreasing involution in \eqref{pack:rhombus} interchanges $\alpha$ and
$2\alpha$. A degree-three vertex is consequently incident to at most
one triangle; a degree-four vertex to at most three. Every contact
three-cycle is facial, by its empty small region.

\subsection{The exact finite graph obligation}
Let $k$ be the number of degree-three vertices. Edge and Euler counting
show $k$ is even and there are $8-k$ triangular and $2+k/2$
quadrilateral faces. If $k=8$, the cubic quadrangulation is bipartite:
every cycle is, modulo two, a sum of quadrilateral face boundaries.
Its bipartition has four vertices in each part, so it is $K_{4,4}$ minus
a perfect matching, namely the cube.

For $k=2,4,6$, the following finite verification gives a spanning cube.
Write $\beta=2\pi-3\alpha$. At a degree-four vertex a quadrilateral
angle equals $\beta$ if three incident
faces are triangles, and is less than $\beta$ otherwise. At a
degree-three vertex it is greater than $\beta$ if one incident face is
a triangle; with no incident triangle we impose no comparison.
These comparisons follow by summing the other incident angles and using
\eqref{pack:angles}. Opposite corners cannot have unequal determined classes.

Here is the complete enumeration obligation and procedure. Label the
$k$ degree-three vertices first. At vertex $i$, its edges to smaller
vertices are already fixed. Choose every subset of the requisite size
from larger vertices with positive degree remainder. Reject only impossible
remainders, and restore all data on backtracking. Enumerate all three-cliques;
require exactly $8-k$, with the incidence bounds just proved, and reject
$K_4$. Enumerate all simple four-cycles up to rotation and reversal,
rejecting a contact diagonal and unequal determined opposite classes.
For every edge set its remaining face need equal to two minus its
number of triangular incidences. Recursively choose an edge with positive
need and try every unused eligible four-cycle through it. Decrease its
four edge needs, recurse, and restore them. A branch is rejected if
fewer eligible four-cycles remain than its required need. Sorted sets
of chosen four-cycle indices remove only repeated search orders.

The extended checker additionally requires the necessary sphere-map
conditions: edge use twice, Euler characteristic two, connected cyclic
vertex links, compatible face orientations, and connected face adjacency.
Every actual embedded contact graph passes these tests. No isomorphism
pruning is used. For each surviving graph, delete in turn every matching
between degree-four vertices that covers them exactly once. Verify that
one deletion leaves a connected cubic bipartite graph with parts of size
four. This is precisely a spanning cube. The search is complete, since
a spanning cubic subgraph deletes no edge at a degree-three vertex and
exactly one edge at each degree-four vertex.

{\raggedright
The programs \path{tail/enumerate_contact_maps_extended.py} and
\path{tail/verify_packing.py} implement these exact
predicates and reject any surviving graph without a spanning cube.
Their fresh replay succeeds. The enumeration is finite and uses integer
combinatorics; no numerical geometric embedding is used as a premise.\par}

\subsection{Four-regular configurations}
For $k=0$, Euler counting gives eight triangles and two quadrilaterals.
Every vertex needs a nontriangular face. The two quadrilaterals are
therefore disjoint and cover all vertices. Every corner has angle
$\beta=2\pi-3\alpha$. Their cyclic Gram matrices have adjacent entries
$t$ and a common diagonal entry $r$. Eigenvalues are $1-r$ twice and
$1+r\pm2t$. Positivity, rank at most three, and distinctness give $r=2t-1$.
Consequently
\[
 \cos\beta=-\frac{1-t}{1+t}=\cos(3\alpha),\qquad
 1-2t-7t^2=0.
\]
Thus $t=a$. Each quadrilateral is a Euclidean square of center norm
$\sqrt a$ and radius $\sqrt{1-a}$, by its Gram matrix. Place one square
at height $h=\sqrt a$. A vertex contacting the two upper corners
$(r,0,h),(0,r,h)$ has the form $(x,x,z)$, where $x=(a-hz)/r$.
Its possible heights solve
\[
 (1+a)z^2-4ahz+2a^2+a-1=0,
\]
and are $-h$ and $h(1+5a)/(1+a)$. The second violates the cutoff with
the opposite upper corner. The four edges of the upper square therefore
force the four lower vertices at height $-h$ and offset $\pi/4$, using
$2a/(1-a)=1/\sqrt2$. This is $P$.

\subsection{Cube angles and a realized packing deformation}
\begin{lemma}[The inherited cube faces]\label{r3:topology-cube-faces}
Suppose eight distinct spherical points have maximum pairwise inner
product $t\in(0,3/10)$, and their noncrossing short contact arcs
contain an abstract spanning cube. Suppose also that the small convex
region of every contact four-cycle contains no other packing point.
Then the inherited cube embedding has six faces, each the small convex
region of a contact four-cycle. Every contact edge omitted from the
cube is a diagonal of one of these facial quadrilaterals.
\end{lemma}
\begin{proof}
The cube is connected and two-connected, so each boundary of a face
of its sphere embedding is a simple cycle. Euler's formula gives six
faces. As the cube is bipartite, every facial cycle has length at
least four; the total of their lengths is twice its twelve edges,
namely twenty-four. Thus all six boundaries are four-cycles.

For each such cycle, the small convex disk contains none of the
remaining four cube vertices, by hypothesis. Its complementary disk
contains them all and consequently cannot be a face. The facial disk
must therefore be the small disk. An omitted contact arc has no
interior cube vertex and cannot cross or overlap any cube edge. Its
open arc lies in one face, and its two endpoints lie on that face's
boundary. They cannot be adjacent, since the unique short arc between
adjacent corners is already a cube edge. They are therefore opposite
corners, as claimed.
\end{proof}

The empty-quadrilateral property proved above verifies the hypothesis
of Lemma~\ref{r3:topology-cube-faces}. Label the vertices of a spanning cube $(x,y,z)\in\{\pm1\}^3$. On the face
$x=\sigma$, express the angle as $m_{1,\sigma}+e_{1,\sigma}yz$, and use
cyclic analogues on the other faces. The eight vertex angle sums are
\[
 m_{1,x}+e_{1,x}yz+m_{2,y}+e_{2,y}xz+m_{3,z}+e_{3,z}xy=2\pi.
\]
Expanding in the eight independent sign monomials gives
\[
 m_{i,+}=m_{i,-}=M_i,\quad e_{i,+}=-e_{i,-}=E_i,\quad
 \sum_iM_i=2\pi,\quad\sum_iE_i=0.
\]
Set $c=(1-t)/(1+t)$ and $L=\alpha/2$. The rhombus identity is equivalent to
\[
 M_i=m_t(E_i)=\arccos(-c\cos E_i),\qquad |E_i|\le L.
\]
The weak endpoints allow added contact diagonals. All angle branches are
in $(0,\pi)$, and
\[
 m_t''(E)=-\frac{c(1-c^2)\cos E}{(1-c^2\cos^2E)^{3/2}}<0.
\]
The plane $\sum E_i=0$ cuts the cube $|E_i|\le L$ in the hexagon with
vertices the permutations of $(L,-L,0)$. At these vertices the common
sum of means is
\[
 3\alpha(t)+\beta_{\square}(t),\qquad
 \beta_{\square}(t)=2\operatorname{arccot}\sqrt t.
\]
Both summands decrease strictly with $t$, and their sum equals $2\pi$
exactly at $t=a$. For $t\le a$, strict concavity therefore forces $t=a$
and a hexagon vertex. The four faces with angles $\alpha,2\alpha$
then supply four distinct contact diagonals forming a perfect matching:
at each cube vertex exactly one of the first two face families has angle
$2\alpha$. Every degree becomes four, and the preceding classification
forces $P$.

For the stronger local assertion allow $t<3/10$ and rescale $E_i=L(t)z_i$.
The fixed domain is $K=\{z\in[-1,1]^3:\sum z_i=0\}$, and the realization
equation is
\[
 F(t,z)=\sum_i m_t(L(t)z_i)-2\pi=0.
\]
This is strictly concave in $z$ and strictly decreasing in $t$. Indeed,
with $C=\sqrt{1+2t}$,
\[
 \alpha'=-\frac1{(1+t)C},\quad c'=-\frac2{(1+t)^2},\quad
 \frac{d}{dt}m_t(E)=\frac{c'\cos E-c\sin E\,E'}{\sin m_t(E)}.
\]
For $E=z\alpha/2$, the ratio of the possibly positive second numerator
term to the magnitude of the first is at most
\[
 \frac{(1-t)|z|\tan|E|}{4C}
 \le\frac{1-t}{4(1+2t)}<1.
\]
If $z$ is not a vertex of $K$, choose a line segment in $K$ with $z$ in
its relative interior. Strict concavity makes $F(t,z')<0$ for points
arbitrarily near $z$ on at least one side. The implicit-function theorem
then gives nearby $t'<t$ with $F(t',z')=0$. At a vertex of $K$, the equation
instead forces $t=a$ and the packing antiprism as above.

\begin{lemma}[Spherical gluing and continuous realization]
\label{r3:topology-spherical-gluing}
Attach six convex spherical quadrilaterals according to the faces of
an oriented abstract cube, identifying matching edges by arclength.
Assume that each edge length is in $(0,\pi)$, each face lies in an
open hemisphere and has all angles in $(0,\pi)$, matching edge
lengths agree, and the three face angles at each cube vertex sum to
$2\pi$. Then the glued surface admits an isometry to the unit sphere.
In particular its vertices and edges form an embedded realization of
the cube. If the face shapes vary continuously near a given such
gluing, the realization can be chosen continuously after fixing the
position and orientation of one face.
\end{lemma}
\begin{proof}
The quotient is a compact connected topological sphere, since its
gluing is the boundary complex of a cube. Face interiors already have
spherical charts. At an interior edge point, place its two incident
faces on opposite sides of their common great circle; this supplies
a spherical chart across the edge.

At a vertex choose a radius smaller than the distances to every
other vertex and every nonincident edge in each of its three faces.
Each incident piece is then a spherical sector with metric
\[
 dr^2+\sin^2r\,d\phi^2.
\]
The three sectors glue cyclically, and their angular widths sum to
$2\pi$. Together they are exactly a spherical disk with this polar
metric. This supplies the vertex chart, including the center, and
shows explicitly that no cone singularity or residual vertex holonomy
remains. All chart transitions are restrictions of orientation
preserving isometries of the unit sphere.

Start with an isometry of one chart into the unit sphere and continue
it along paths by these transitions. Continuation exists along every
path by compactness of its parameter interval. It is independent of
path: a homotopy of paths can be subdivided into sufficiently small
rectangles lying in charts, on each of which the transition cocycle
telescopes. The source sphere is simply connected, so all paths with
the same endpoints are homotopic. We obtain a global local isometry
$f$ into the unit sphere.

Its image is open because $f$ is a local homeomorphism and closed
because the source is compact, hence is the entire unit sphere. A
proper local homeomorphism is a covering: each fiber is finite, and
disjoint chart neighborhoods of its points can be shrunk to a common
image neighborhood avoiding the image of their compact complement.
A connected covering of the simply connected sphere has one sheet.
Consequently $f$ is a bijective local isometry and thus an isometry.

Finally, choose a fixed spanning tree of the six face adjacencies.
Position the first face continuously, and position successive faces
across their shared edges on the prescribed opposite side. All these
operations are continuous for the present nondegenerate short edges
and face angles. The resulting face maps agree with the unique global
development just constructed. In particular all eight developed
vertices vary continuously and agree with the original realization
at the original parameter value.
\end{proof}

We apply this lemma to realize the angle variations as nearby configurations. A prescribed rhombus with diagonal inner products $r,s$ has
corners
\[
 h_1m\pm\sqrt{1-h_1^2}\,e_1,\qquad
 h_2m\pm\sqrt{1-h_2^2}\,e_2,
\]
where $h_1=\sqrt{(1+r)/2}$, $h_2=\sqrt{(1+s)/2}$, $h_1h_2=t'$, and
$(m,e_1,e_2)$ is orthonormal. The cosine law recovers $r,s$ continuously
from the angles. These six convex spherical rhombi have matching short
edge lengths and vertex angle sums $2\pi$. Lemma~\ref{r3:topology-spherical-gluing}
therefore supplies an embedded realization varying continuously from the
original one.
The bounds $|z_i|\le1$ keep all facial diagonals at inner product at most
$t'$. Every original contact outside the cube is such a diagonal, because
its short arc cannot cross or overlap a cube edge or pass through another
vertex. All remaining strict noncontacts keep their slack by continuity.
We have constructed a genuine smaller cutoff, contradicting local optimality.

This proves the first assertion of Theorem~\ref{r3:packingclassification}.
For the global assertion, compactness gives a global packing minimizer;
choose one with fewest contacts. Since $P$ attains $a<3/10$, the first
assertion forces this chosen minimizer to be $P$. It has sixteen contacts.
Every other global minimizer has at least sixteen by the selection and
at most sixteen by the degree bound, so it is four-regular and the preceding
geometric argument gives the same congruence class. This completes the
proof of the packing theorem and its equality case.


\section{Quantitative geometry near the packing antiprism}\label{sec:geometry}
Throughout this section let
\[
\begin{gathered}
 a=\frac{2\sqrt2-1}{7},\quad P=X(a),\quad \tau=2(1-a),\\
 \kappa=1+\frac1{\sqrt2},\quad W=8(\kappa+1),\quad d_w=W/4.
\end{gathered}
\]
The contacts of $P$ are its eight square edges and eight nearest cross
edges. Their squared distance is $\tau$; every other squared distance
is at least $2\tau$. Give square contacts weight $\kappa$ and cross
contacts weight $1$. Write $\mathcal C$ for this weighted contact set.

We use the tangent chart and rotation slice
\[
 y_i(q)=\sqrt{1-|q_i|^2}\,P_i+q_i,\qquad
 V=\{(q_i):q_i\perp P_i,\ \sum_iP_i\times q_i=0\}.
\]
Put $r=|q|$, $f_{ij}(q)=|y_i(q)-y_j(q)|^2$, $t_i=|q_i|^2$ and
$t_e=t_i+t_j$ for $e=\{i,j\}$. The linear and quadratic contact terms are
\[
 \ell_e(q)=-2(P_i\cdot q_j+P_j\cdot q_i),\qquad
 Q_e(q)=a(t_i+t_j)-2q_i\cdot q_j.
\]
Let $R:q\mapsto(\ell_e(q))_e$ on the full sixteen-dimensional tangent
space, and $D_w=\operatorname{diag}(w_e)$. Direct substitution gives the
positive stress identity $R^Tw=0$. Each vertex has weighted degree $d_w$.

\subsection{Contact matrices and chart remainders}
The constants needed below are
\begin{equation}\label{geo:constants}
 \rho=\frac3{10},\qquad \gamma=\frac32,\qquad h_0=\frac{11}{50},
 \qquad \mu=\frac{59546}{15625},\qquad \delta=\frac{289}{10000}.
\end{equation}
We use the unnormalized weighted norm
$\|x\|_w^2=\sum_ew_ex_e^2$ and write
$C_wx=x-W^{-1}(\sum_ew_ex_e)\mathbf1$ for weighted centering.

We construct the coordinates for the matrix estimates as follows.
With $\theta_j=j\pi/4$ and $\sigma_j=(-1)^j$, use the tangent frames
\[
 u_j=(-\sin\theta_j,\cos\theta_j,0),\qquad
 v_j=(-\sqrt a\cos\theta_j,-\sqrt a\sin\theta_j,
                           \sigma_j\sqrt{1-a}).
\]
The three rotation columns in these frames give their orthogonal
projection by a $3$-by-$3$ Gram inverse. Projecting the coordinate
vectors onto $V$ and applying Gram--Schmidt to thirteen independent
columns gives an orthonormal basis. The column choices are witness
data; the checker proves that every normalization denominator is
positive. Orthogonality follows from the defining exact operations.
From now on, $R$, $M_e$ and $T_e$ denote the matrices in this basis,
so that $\ell_e(q)=R_eq$, $Q_e(q)=q^TM_eq$ and $t_e=q^TT_eq$.
All entries are defined by radicals and finite algebraic operations,
with outward rational enclosures constructed by
\path{tail/exact_geometry.py}.

\begin{lemma}[Contact matrix bounds]\label{geo:spectrum}
The kernel of the contact differential on the full tangent space is
exactly the space of infinitesimal rotations. For every $q\in V$,
\[
 \|Rq\|_2^2\ge3|q|^2,\qquad \|Rq\|_w^2\ge\mu|q|^2.
\]
Moreover, with $\bar M=W^{-1}\sum_ew_eM_e$,
\begin{equation}\label{geo:variancegram}
 \sum_ew_e(M_e-\bar M)^2\preceq4I.
\end{equation}
\end{lemma}
\begin{proof}
{\raggedright
The program \path{tail/verify_variance.py} constructs
these matrices from the preceding coordinates. Outward $LDL^T$
elimination gives thirteen positive pivots for each of\par}
\[
 R^TR-3I,\qquad R^TD_wR-\mu I,\qquad
 4I-\sum_ew_e(M_e-\bar M)^2.
\]
These finite checks prove the claimed comparisons on $V$.
The three independent rotation vectors are annihilated by the
contact differential, while its restriction to the thirteen-dimensional
complement is injective. This identifies its kernel.
\end{proof}

For $r\le\rho$, expansion of the chart gives
\begin{equation}\label{geo:remainders}
 |f_e-\tau-\ell_e|\le\gamma t_e,
 \qquad |f_e-\tau-\ell_e-Q_e|\le h_0t_e.
\end{equation}
Indeed, setting $h_i=\sqrt{1-t_i}$, the full remainder equals
\[
 2a(1-h_ih_j)+2(1-h_i)P_i\cdot q_j
 +2(1-h_j)P_j\cdot q_i-2q_i\cdot q_j.
\]
Its absolute value divided by $t_e$ is at most
\[
 1+\frac{2a+\sqrt{2(1-a^2)}\rho}{1+\sqrt{1-\rho^2}}<\gamma.
\]
After subtracting $Q_e$, the first term becomes $a(h_i-h_j)^2$.
The inequality $|P_i\cdot q_j|\le\sqrt{1-a^2}|q_j|$ then bounds the
remaining absolute error divided by $t_e$ by
\[
 \frac{a\rho^2}{4(1-\rho^2)}+
 \frac{\sqrt{2(1-a^2)}\rho}{1+\sqrt{1-\rho^2}}<h_0.
\]
The scalar checker verifies both strict comparisons.

\subsection{A quartic certificate for global packing separation}
\begin{lemma}[Global separation]\label{geo:gap}
If the distance to the congruence orbit of $P$, measured in the
product Euclidean norm, is at least $\rho$, then
\[
 \min_{i<j}|y_i-y_j|^2<\tau-\delta.
\]
\end{lemma}
\begin{proof}
First consider a configuration on the boundary of this orbit ball
and choose a closest alignment with $P$. Put $r_i=|y_i-P_i|$.
Then $P_i\cdot y_i=1-r_i^2/2>0$, so its tangent projections $q_i$
are valid chart coordinates. Differentiating the alignment distance
under rotations gives $q\in V$, and
\[
 |q|^2=\rho^2-\frac14\sum_i r_i^4,
 \qquad \rho^2(1-\rho^2/4)\le|q|^2\le\rho^2.
\]
Writing $q=\rho u$ in the orthonormal coordinates above, define
\[
 \begin{aligned}
 g_e(u)&=\delta+\rho R_eu+\rho^2u^T(M_e+h_0T_e)u
                       &&(e\in\mathcal C),\\
 g_{17}(u)&=1-|u|^2,\qquad
 g_{18}(u)=|u|^2-\frac{391}{400}.
 \end{aligned}
\]
If all contact distances were at least $\tau-\delta$, the upper
remainder bound in \eqref{geo:remainders} would imply all eighteen
inequalities $g_j(u)\ge0$.

The exact certificate checked by
\path{tail/verify_sos_gap.py} is the polynomial identity
\begin{equation}\label{geo:sosidentity}
 -1=z_2(u)^TQ_0z_2(u)
       +\sum_{j=1}^{18}z_1(u)^TQ_jz_1(u)g_j(u),
 \qquad Q_j\succ0.
\end{equation}
Here $z_1$ and $z_2$ list all monomials of degree at most one and two,
respectively, so the Gram dimensions are $14$ and $105$.
For precision, the checker keeps the eighteen rational multiplier
matrices fixed and corrects a proposed free matrix $Q_0^{\rm rat}$ as follows.
Let $r_\alpha$ be a coefficient residual and $n_\alpha$ the number
of ordered Gram entries representing $u^\alpha$. Every monomial of
degree at most four occurs, so $n_\alpha>0$. Add
$r_\alpha/n_\alpha$ to each corresponding entry of $Q_0^{\rm rat}$.
This makes all $2380$ coefficients exact and preserves symmetry.
If $e$ bounds the absolute value of every added entry, then the
correction $\Delta$ satisfies $\|\Delta\|_2\le105e$.
A rational congruence matrix $A_0$ is checked to satisfy
\[
 \|A_0Q_0^{\rm rat}A_0^T-I\|_2+
        105e\,\|A_0\|_1\|A_0\|_\infty<1.
\]
Thus the corrected matrix is positive definite. The other eighteen
matrices have congruence defects below one without correction.
Integer matrix products and rational interval bounds prove these
comparisons; numerical eigenvalues are not proof inputs.
Identity \eqref{geo:sosidentity} now contradicts $g_j(u)\ge0$.

To pass from the boundary to the entire exterior, minimize $T(Y)$
on the compact exterior and choose a minimizer with fewest contacts.
An interior minimizer with $T<3/10$ would be contact-minimal locally
optimal, so Theorem~\ref{r3:packingclassification} would make it
congruent to $P$, which is impossible. If $T\ge3/10$, its distance
deficit is at least $2(3/10-a)>\delta$, an exact checked comparison.
The remaining boundary case was just excluded. This minimizer has
the largest minimum squared distance in the exterior, proving the
claim for every exterior configuration.
\end{proof}

\subsection{A bound for the signed contact mean}
Let $K=ad_wI-\operatorname{Adj}_w$ be the weighted graph stress matrix.
Then $KP=0$. For $\Delta Y=Y-P$, the unit lengths give the exact identity
\[
 S(q):=\sum_ew_e(f_e(q)-\tau)
       =\sum_{j=1}^3(\Delta Y^{(j)})^TK\Delta Y^{(j)}.
\]
In the tangent and radial frames write this as
\begin{equation}\label{geo:meanblocks}
 S(q)=q^THq+2q^TBn+n^TDn,\qquad n_i=\sqrt{1-t_i}-1.
\end{equation}
The letters $H,B,D$ in this subsection denote these blocks only.

\begin{lemma}\label{geo:mean}
For $q\in V$ and $r=|q|<1$,
\[
 S(q)\le\left[\frac{413}{100}+\frac34r^2+
 \frac{(4r+231/100)r^2}{(1+\sqrt{1-r^2})^2}\right]r^2.
\]
\end{lemma}
\begin{proof}
We first specify the small exact matrix certificate used in the proof.
Put $\alpha=413/100$. The cyclic Fourier decomposition of
\eqref{geo:meanblocks} gives, for $\omega=k\pi/4$, the rephased blocks
\[
 H_k=\begin{pmatrix}t_k&b_k\\b_k&m_k\end{pmatrix},\quad
 B_k=\binom{-e_k}{f_k},\quad D_k=d_k,
\]
where, with $h=\sqrt a$,
\begin{align*}
 t_k&=\sqrt2(1-\cos\omega),\\
 m_k&=\sqrt2+[2-(2+\sqrt2)a]\cos\omega-2(1+a)\cos(2\omega),\\
 b_k&=2h[(\sqrt2/2)\sin\omega+\kappa\sin(2\omega)],\\
 e_k&=2\sqrt{1-a}[(\sqrt2/2)\sin\omega+\kappa\sin(2\omega)],\\
 f_k&=2h\sqrt{1-a}[(1+\sqrt2/2)\cos\omega-\kappa\cos(2\omega)],\\
 d_k&=2a[\kappa+1-\cos\omega-\kappa\cos(2\omega)].
\end{align*}
The rephasing multiplies the azimuthal Fourier coordinate by $i$.
For $k=0,4$, the two coordinate vectors are diagonal modes with
entries $\cos(k\theta_j)/\sqrt8$. For $k=1,2,3$, the eigenvalues are
\[
 h_k^\pm=\tfrac12(t_k+m_k\pm\sqrt{(t_k-m_k)^2+4b_k^2}).
\]
Normalize $(b_k,h_k^\pm-t_k)$ to a unit vector $(\xi,\zeta)$.
The corresponding real modes have tangent pairs
$(-\xi\sin(k\theta_j),\zeta\cos(k\theta_j))/2$ and
$(\xi\cos(k\theta_j),\zeta\sin(k\theta_j))/2$.
Trigonometric orthogonality and the orthogonality of the two eigenvectors
of each symmetric block give an orthonormal eigenbasis $V_0$ of $H$.
All displayed normalizers are nonzero, as checked by exact intervals.

Write $q=V_0x$, $V_0^THV_0=\operatorname{diag}(h_i)$ and
$\delta_i=\alpha-h_i$. The certificate checks $\delta_i>0$,
$e_k^2+f_k^2<16$ and $d_k<231/100$. Thus
$\|B\|_2<4$ and $D\preceq(231/100)I$.
Let $C_3(q)=-q^TB(t_i)_i$, and write its fully symmetric coefficient
tensor in the $x$ coordinates as $C_3(V_0x)=\sum_{ijk}T_{ijk}x_ix_jx_k$.
Define
\[
 U_{ijk}=\frac{3T_{ijk}}{\delta_i+\delta_j+\delta_k},\qquad
 N_{i,(j,k)}=\sqrt{\delta_i}\,U_{ijk}.
\]
The ordered pair $(j,k)$ indexes the 256 columns of $N$. The final
certificate is the positive definiteness of $3I-NN^T$, proved by
sixteen positive outward $LDL^T$ pivots. Every entry is reconstructed
from the formulas above by
\path{tail/verify_wholeball_mean.py}.
This is a finite algebraic comparison; a numerical eigenvalue is not
used as a bound.

Symmetry of $U$ gives
$C_3(V_0x)=\sum_{ijk}\delta_iU_{ijk}x_ix_jx_k$.
With $\Delta=\sum_i\delta_ix_i^2=q^T(\alpha I-H)q$, Cauchy--Schwarz
therefore yields $|C_3(q)|\le\sqrt{3\Delta}\,r^2$.
The radial identity, including its entire remainder, is
\[
 n_i=-\frac{t_i}{2}+e_i',\qquad
 e_i'=-\frac{t_i^2}{2(1+\sqrt{1-t_i})^2}.
\]
For $L=1+\sqrt{1-r^2}$, this gives
$|e'|\le r^4/(2L^2)$ and $|n|^2\le r^4/L^2$.
Substitution into \eqref{geo:meanblocks} proves
\[
 S(q)\le\alpha r^2-\Delta+\sqrt{3\Delta}\,r^2
                       +4r^5/L^2+(231/100)r^4/L^2.
\]
Finally $-\Delta+\sqrt{3\Delta}\,r^2\le3r^4/4$ by completing a
square. This proves the bound on the whole ball.
\end{proof}

\subsection{Two bounds for the centered contact norm}
The matrix bound \eqref{geo:variancegram} and Cauchy--Schwarz give
\[
 \|C_wQ(q)\|_w^2
 =\sum_ew_e\bigl(q^T(M_e-\bar M)q\bigr)^2\le4r^4.
\]
For the higher error $h=f-\tau-Rq-Q(q)$, the bounds
$t_e\le r^2$ and $\sum_ew_et_e=d_wr^2$ imply
$\|h\|_w\le h_0\sqrt{d_w}\,r^2$.
Centering is an orthogonal projection in the weighted norm. Hence
\begin{equation}\label{geo:sigma}
 \|C_w(f-\tau)-Rq\|_w\le\sigma r^2,
 \qquad \sigma=2+h_0\sqrt{d_w}.
\end{equation}
A stronger lower bound retains the interaction of the linear and
quadratic contact terms.

\begin{lemma}[Cubic absorption]\label{geo:newvariance}
Put $\mu_0=37/10$, $B_0=528/25$ and $\epsilon_0=10^{-7}$. For
$q\in V$ with $r=|q|\le\rho$,
\begin{equation}\label{geo:variancelower}
 \|C_w(f-\tau)\|_w\ge\mathcal F(r),\qquad
 \mathcal F(r)=r\sqrt{\mu_0-\tfrac14B_0r^2-\epsilon_0r}
                         -h_0\sqrt{d_w}\,r^2.
\end{equation}
\end{lemma}
\begin{proof}
Let $G=R^TD_wR$ and $A=G-\mu_0I$. In addition to the comparisons in
Lemma~\ref{geo:spectrum}, the exact variance checker proves
\[
 A\succ0,\qquad UU^T\preceq B_0A^{-1},
\]
for the rational symmetric tensor $U$ in
\path{tail/variance_witness.json}, viewed as a
$13$-by-$169$ matrix with ordered tensor coordinates. It also proves
\[
 \left|2\sum_ew_e\ell_e(q)Q_e(q)
       -q^TAU(q\otimes q)\right|\le\epsilon_0r^3.
\]
The matrix comparisons use thirteen positive outward $LDL^T$ pivots.
For the last bound, the checker forms every exact cubic coefficient
residual and bounds their absolute sum by $\epsilon_0$; each cubic
monomial has absolute value at most $r^3$.

Since $R^Tw=0$, centering fixes $Rq$. Dropping the nonnegative
quadratic-square term gives
\[
 \|C_w(Rq+Q(q))\|_w^2
 \ge q^TGq+2\sum_ew_e\ell_e(q)Q_e(q).
\]
Write $\Delta=q^TAq$. The matrix comparison and
$|q\otimes q|=r^2$ give
$|q^TAU(q\otimes q)|\le\sqrt{B_0\Delta}\,r^2$. Therefore
\[
 \|C_w(Rq+Q(q))\|_w^2
 \ge\mu_0r^2+\Delta-\sqrt{B_0\Delta}\,r^2-\epsilon_0r^3
 \ge\mu_0r^2-\tfrac14B_0r^4-\epsilon_0r^3.
\]
The last inequality completes a square. The radicand is positive
on $[0,\rho]$ after its factor $r^2$ is removed, as checked at
$\rho$. Finally subtract the centered higher-error norm
$h_0\sqrt{d_w}\,r^2$ by the triangle inequality.
\end{proof}

\section{The unbounded exponent range}\label{sec:tail}
We prove Theorem~\ref{thm:main} for $s\ge124$. Set $p=s/2$ and
$p_0=62$. Constants concerning the packing chart retain the meaning
fixed in Section~\ref{sec:geometry}.
All scalar comparisons in this section are checked by
\path{tail/verify_tail124.py}, using outward rational
intervals and explicit exponential and logarithmic remainder bounds.
Together with the monotonicity arguments below, these comparisons apply
to every real $p\ge p_0$.

\subsection{A necessary condition for stationarity}
For a stationary configuration and $f_{ij}=|y_i-y_j|^2$, the tangential
force equations imply
\begin{equation}\label{r5:centroid}
 \sum_{i<j}f_{ij}^{-p}(y_i+y_j)=0.
\end{equation}
Indeed, the force equation at $i$ can be written
\[
 \sum_{j\ne i}f_{ij}^{-p-1}(y_j-y_i)
      =-\frac12\sum_{j\ne i}f_{ij}^{-p}y_i.
\]
Summing in $i$ cancels the left side and gives \eqref{r5:centroid}.

\begin{lemma}[Force efficiency]\label{r6:force}
For $N\ge3$ distinct stationary points on the unit two-sphere, choose a
closest pair $\{i,k\}$, put $m=f_{ik}$, $D=m^{-p}$, and
$\ell=\sqrt{4-m}$. Let $U$ be the sum of $(m/f)^p$ on the $2(N-2)$ other
edges incident to either endpoint. Then
\begin{equation}\label{r6:forcebound}
 U\ge2,\qquad \frac ED\ge6-m+A(m)(U-2),\qquad
 A(m)=\frac{(\ell-1)^2+2/\ell}{4}>0.
\end{equation}
For $N>3$ one has $U>2$ and $E/D>6-m$.
\end{lemma}
\begin{proof}
The identity $\sum_{a<b}|y_a-y_b|^2=N^2-|\sum_a y_a|^2$ implies
$0<m\le2N/(N-1)\le3$, so the closest pair is not antipodal. On an adjacent
edge of squared length $f\ge m$, set
\[
 x=(m/f)^p,\qquad q(f)=\sqrt{\frac{m(4-f)}{f(4-m)}}\in[0,1].
\]
The tangent projection of the other vertex, divided by $f$, has norm
$\sqrt{(4-f)/f}/2$. Isolating the closest-edge term in the endpoint
force equation and taking norms gives $1\le\sum xq$ at each endpoint.
Consequently $\sum_{\rm adj}xq\ge2$ and $U\ge2$.

An exact pointwise estimate is
\[
 f-m\le K(m)(1-q),\qquad K(m)=(4-m)/2+\sqrt{4-m}.
\]
Indeed $f-m=(1-q)(4-m)f(1+q)/4$, while Cauchy--Schwarz and
$(f-2)^2+f(4-f)=4$ give
\[
 f(1+q)=f+\sqrt{\frac m{4-m}}\sqrt{f(4-f)}\le2+\frac4{\sqrt{4-m}}.
\]
These formulas include the endpoints without dividing by $1-q$. Hence
$M:=\sum_{\rm adj}x(f-m)\le K(m)(U-2)$.

The energy-weighted centroid identity \eqref{r5:centroid} holds for every
stationary configuration. Project it onto $u=(y_i+y_k)/\ell$. An adjacent
edge has projection at least
\[
 h_m(f)=\ell(1-f/4)-\frac{\sqrt m}{2}\sqrt{f(1-f/4)},
\]
as follows by resolving $u$ and the other vertex along and perpendicular
to its incident endpoint. Every disjoint edge has projection at least
$-\ell$. If $V_0$ denotes the normalized energy on disjoint edges, then
\[
 V_0\ge1+\ell^{-1}\sum_{\rm adj}x h_m(f),\qquad E/D=1+U+V_0.
\]
The function $h_m$ is convex on $[m,4]$ and
\[
 h_m(m)/\ell=1-m/2,\qquad
 h_m'(m)/\ell=-\frac{3-m}{2(4-m)}\le0.
\]
Its tangent lower bound, followed by the upper bound on $M$, yields
\[
 \frac ED\ge6-m+(U-2)\left(2-\frac m2-
                  \frac{(3-m)K(m)}{2(4-m)}\right).
\]
The bracket equals $(5-m)/4-(3-m)/(2\ell)=A(m)$.
If $N>3$ and $U=2$, positivity of every $x$ and $q\le1$ force $q=1$ on
every adjacent edge. Then every such $f=m$, so $U=2(N-2)>2$, a
contradiction. This also proves the strict conclusion.
\end{proof}

For eight points, Lemma~\ref{r6:force} gives the strict consequence
\begin{equation}\label{tail:force}
 E_sm^p>6-m
\end{equation}
for every collision-free stationary configuration.

\subsection{Trial energy and entry into the chart}
Define
\[
\begin{gathered}
 d=\frac{\log\kappa}{\kappa+1},\qquad \lambda=e^{-\kappa d},\qquad
 C_{\min}=\lambda W,\\ D=\frac{(\kappa d)^2}{2},\qquad
 C=\frac{9669627}{625000}.
\end{gathered}
\]
The trial antiprism of height $a+\tau d/(2p)$ has short squared
distances $\tau(1-d/p)$ and $\tau(1+\kappa d/p)$. The inequalities
\[
 -\log(1-x)\le x+\frac{x^2}{2(1-x)},\qquad
 \log(1+x)\ge x-x^2/2
\]
and the endpoint comparison $d^2/[2(1-d/p_0)]<D$ bound the sum of
the sixteen short-pair energies, after multiplication by $\tau^p$,
by $C_{\min}e^{D/p}$.
The other distances contribute less than $16\,2^{-p}<10^{-6}$:
the four opposite-square terms are bounded using
$e^{d+D/p_0}<2$, and the remaining eight squared distances exceed
$2\tau$. The trial height belongs to $(0,1)$.
The exact endpoint comparison therefore gives, for every $p\ge p_0$,
\begin{equation}\label{tail:ceiling}
 \tau^p E_s(X(a_s))\le
 \tau^p E_s(X(a+\tau d/(2p)))
 <C_{\min}e^{D/p}+10^{-6}<C.
\end{equation}
All adverse bounds here decrease with $p$.

Set $m_0=\tau-\delta$. The scalar checker verifies
\[
 0<m_0<\tau,\qquad (6-m_0)(\tau/m_0)^{p_0}>C.
\]
For fixed $m_0$ the left side increases with $p$, and for fixed $p$
the expression $(6-m)(\tau/m)^p$ strictly decreases with $m$ on
$(0,m_0]$. Lemma~\ref{r6:force} therefore implies that every
stationary configuration below the ceiling has $m>m_0$.
Lemma~\ref{geo:gap} places it inside the orbit ball of radius $\rho$.
In particular, this applies both to every possible global minimizer
and to $X(a_s)$, which is fully stationary by the symmetry argument
in Section~1. After closest alignment, both have chart representatives
in $V$ with $r<\rho$.

\subsection{A finite sequence of radius improvements}
We first record the elementary extremal inequality used for the
exponential contact energy.

\begin{lemma}[Two-value exponential extremum]\label{r3:twovalue}
Let positive weights $w_i$ have sum $W$, and suppose
$\sum_iw_iu_i=0$, $v=\sum_iw_iu_i^2>0$. For every subset weight
$b=\sum_{i\in I}w_i$ with $0<b<W$, put
\[
 J_b(v)=b\exp\left(-\sqrt{\frac{v(W-b)}{Wb}}\right)
 +(W-b)\exp\left(\sqrt{\frac{vb}{W(W-b)}}\right).
\]
Then $\sum_iw_ie^{-u_i}\ge\min_b J_b(v)$. Moreover every $J_b$ is strictly
increasing in $v>0$.
\end{lemma}
\begin{proof}
At fixed positive variance and zero mean the feasible set is a compact
ellipsoid section. The two constraint gradients are independent, since
dependence would make all $u_i$ equal and hence the variance zero.
Lagrange multipliers at a minimum give $e^{-u_i}=\alpha+\beta u_i$.
A strictly convex exponential intersects a line in at most two points.
There must be two, with one positive and one negative. If the positive
entries have total weight $b$, the mean and variance determine their
values as $\sqrt{v(W-b)/(Wb)}$ and $-\sqrt{vb/[W(W-b)]}$. This gives the
listed alternatives. Write $v=t^2$. The derivative of $J_b(t^2)$ is a
positive constant times $e^{ct}-e^{-at}>0$, where
$a=\sqrt{(W-b)/(Wb)}$ and $c=\sqrt{b/[W(W-b)]}$. This proves monotonicity.
\end{proof}


Let $v=\partial_bX(b)|_{b=a}$ in the aligned tangent coordinates.
Then $v\in V$, $|v|^2=2/[a(1-a)]$, and its contact differentials are
$-2$ on square contacts and $2\kappa$ on cross contacts. Set
\[
 q_*=(\tau d/(2p))v,\qquad N=\frac d{\sqrt{2a(1-a)}},\qquad
 z_e^*=(p/\tau)\ell_e(q_*).
\]
Thus $z_e^*=-d$ or $\kappa d$, $p|q_*|=\tau N$, and
$e^{-z_e^*}=\lambda w_e$. Define
\[
 \begin{gathered}
 Z_e=(p/\tau)(f_e-\tau),\qquad z_e=(p/\tau)\ell_e(q),\\
 \bar Z=W^{-1}\sum_ew_eZ_e,\qquad u_e=Z_e-z_e^*-\bar Z.
 \end{gathered}
\]
The list $u$ has weighted mean zero. Lemma~\ref{geo:mean}, evaluated
at $\rho$, gives on the entire chart
\begin{equation}\label{tail:mean}
 \bar Z\le\eta pr^2/\tau,\qquad \eta=1/5.
\end{equation}
The coefficient divided by $W$ is increasing with $r$ and is checked
to be below $\eta$ at $\rho$. Since $\log(1+x)\le x$,
\begin{equation}\label{tail:contactenergy}
 \tau^pE_s(q)\ge\sum_{e\in\mathcal C}e^{-Z_e}
              =\lambda e^{-\bar Z}\sum_ew_ee^{-u_e}.
\end{equation}

\begin{lemma}[The outer annulus]\label{tail:outer}
For every real $p\ge p_0$, a configuration in the chart below the
trial ceiling has $r<6/p$.
\end{lemma}
\begin{proof}
Put
\[
 \begin{gathered}
 Z_* =\|z^*\|_w=\sqrt{8\kappa d^2+8\kappa^2d^2},\\
 a_1=\frac1{\sqrt{W(W-1)}},\qquad
 \mathcal H(r)=a_1\mathcal F(r)-\eta r^2.
 \end{gathered}
\]
Lemma~\ref{geo:newvariance} and the triangle inequality imply
\[
 \|u\|_w\ge v(p,r):=p\mathcal F(r)/\tau-Z_*.
\]
The function $\mathcal F(r)/r$ decreases. Consequently, throughout
$6/p\le r\le\rho$,
\[
 v(p,r)\ge\frac{6\mathcal F(\rho)}{\tau\rho}-Z_*>\sqrt W,
\]
where the final strict inequality is checked exactly.
For any possible subset mass $b$ in Lemma~\ref{r3:twovalue}, keep
only the positive exponential term of $J_b(v^2)$. When $v\ge\sqrt W$,
the function
\[
 (W-b)\exp\left(v\sqrt{\frac b{W(W-b)}}\right)
\]
is nondecreasing in $b$: its logarithmic derivative is nonnegative
precisely when $v\ge2\sqrt{b(W-b)/W}$, and the latter bound follows
from $4b(W-b)\le W^2$. Every nonempty subset has mass at least one.
Thus \eqref{tail:mean}--\eqref{tail:contactenergy} imply
\begin{equation}\label{tail:outerenergy}
 \tau^pE_s(q)\ge\lambda(W-1)
            \exp\left(-a_1Z_*+\frac p\tau\mathcal H(r)\right).
\end{equation}

To control the whole interval in $r$ and the whole half-line in $p$,
write $g(r)=\mu_0-\epsilon_0r-(B_0/4)r^2$. Its derivative and second
derivative are negative, and
\[
 (r\sqrt{g(r)})''=
 \frac{g'(r)}{\sqrt{g(r)}}+
 \frac{r g''(r)}{2\sqrt{g(r)}}-
 \frac{r(g'(r))^2}{4g(r)^{3/2}}<0.
\]
Hence $\mathcal H$ is strictly concave and its minimum on
$[6/p,\rho]$ is at an endpoint. The checker proves
$\mathcal H(\rho)>0$ and
\[
 \lambda(W-1)\exp\left(-a_1Z_*+
        \frac{p_0}{\tau}\mathcal H(r)\right)>C
 \qquad\left(r=\frac6{p_0},\ \rho\right).
\]
The quantity $p\mathcal H(\rho)$ increases. The other endpoint is
\[
 p\mathcal H(6/p)=
 6a_1\sqrt{\mu_0-6\epsilon_0/p-9B_0/p^2}
       -\frac{36(a_1h_0\sqrt{d_w}+\eta)}p,
\]
which also increases. Both endpoint comparisons therefore remain
strict for every $p\ge p_0$. Concavity and
\eqref{tail:outerenergy} exclude every $6/p\le r\le\rho$.
\end{proof}

Four further annuli give the final radius. Lemma~\ref{geo:spectrum}
and \eqref{geo:sigma} imply
\begin{equation}\label{tail:variance}
 \|u\|_w\ge
 \sqrt\mu\left(\frac{pr}\tau-N\right)-\frac{p\sigma r^2}\tau.
\end{equation}
Indeed the centered linear list is $(p/\tau)R(q-q_*)$, and its
distance from the actual centered list is at most $p\sigma r^2/\tau$.
Use the exact sequence
\[
 6\longrightarrow\frac{27}{10}\longrightarrow\frac{13}{10}
       \longrightarrow\frac{43}{50}\longrightarrow\frac34.
\]
For each step $K$ to $L$, set
\[
 v_L=\left[\sqrt\mu(L/\tau-N)-\frac{\sigma L^2}{p_0\tau}\right]^2.
\]
The checker verifies positivity of the bracket and, for all $79$
masses $b=i\kappa+j$ with $0\le i,j\le8$ other than $0,W$,
\begin{equation}\label{tail:annulustest}
 \lambda J_b(v_L)>C\exp\left(\frac{\eta K^2}{p_0\tau}\right).
\end{equation}
These tests exclude each entire annulus for every $p\ge p_0$.
In fact the checked inequality $\sqrt\mu-2\sigma\rho>0$ makes
the right side of \eqref{tail:variance} increase in $r\le\rho$.
For $r\ge L/p$ it is at least $\sqrt{v_L}$, since increasing $p$
decreases the adverse error term. For $r\le K/p$, the mean is at
most $\eta K^2/(p\tau)$. Lemma~\ref{r3:twovalue} and its monotonicity,
together with \eqref{tail:annulustest}, contradict the ceiling throughout
$L/p\le r\le K/p$. It follows that
\begin{equation}\label{tail:finalradius}
 |q|<3/(4p).
\end{equation}
The argument uses only the energy ceiling and chart entry, and hence
applies to both the candidate and every possible global minimizer.

\subsection{Strict convexity on a common containing set}
Let $K_f=3/4$ and define
\[
 \epsilon=\frac{\gamma K_f^2}{p_0\tau},\qquad
 C_*=Ce^\epsilon,\qquad L_p(q)=\sum_{e\in\mathcal C}e^{-z_e},
\]
\begin{equation}\label{tail:convexset}
 \mathcal K_p=\{q\in V:|q|\le K_f/p,\ L_p(q)\le C_*\}.
\end{equation}
The set is convex because each $z_e$ is linear. On the final ball,
\eqref{geo:remainders} gives $|Z_e-z_e|\le\epsilon$ and hence
$\tau^pE_s(q)\ge e^{-\epsilon}L_p(q)$.
Thus the candidate and every global minimizer have representatives
in this same set. The following bounds hold everywhere on it.

For $u_e=z_e-z_e^*$, the weighted mean is zero and
$L_p=\lambda\sum_ew_ee^{-u_e}$.
If one coordinate of weight $w$ has value $u>0$, weighted Jensen on
the remaining coordinates gives
\[
 L_p\ge\lambda\{we^{-u}+(W-w)e^{wu/(W-w)}\}.
\]
The right side increases for $u>0$. The two endpoint tests at $z_e=11/10$,
using $w=\kappa,z_e^*=-d$ or $w=1,z_e^*=\kappa d$, exceed $C_*$.
Therefore $z_e<11/10$ for every contact, while $z_e\ge-\log C_*$ follows
from its own summand. The checked comparisons imply
\begin{equation}\label{tail:contactbounds}
 1-3/p<f_e/\tau<1+3/p,\qquad (\tau/f_e)^{p+2}>3/10.
\end{equation}
For the second inequality use
$(1+(11/10+\epsilon)/p)^{p+2}\le e^{(11/10+\epsilon)(1+2/p_0)}<10/3$.

There is a stronger bound on the incident contact energies.
Put $x_e=e^{-z_e}$, $L=\sum_ex_e$ and $S_i=\sum_{e\ni i}x_e$.
The incident and complementary stress masses are $W/4$ and $3W/4$.
Weighted Jensen for the logarithm in each group, together with
$\sum_ew_e\log x_e=0$, gives
\begin{equation}\label{tail:incidententropy}
 S_i(L-S_i)^3\ge\frac{27}{256}C_{\min}^4.
\end{equation}
Indeed, for a group of stress mass $b$ and total $x$-mass $S$,
\[
 \sum_{e\text{ in group}}w_e\log x_e
 \le b\log(S/b)+\sum_{e\text{ in group}}w_e\log w_e.
\]
Adding the two inequalities and using
$W\log\lambda=-8\kappa\log\kappa$ proves
\eqref{tail:incidententropy} by exponentiation.
Since $L\le C_*$, the checker comparisons
\[
 C_*/4<5<C_*,\qquad 5(C_*-5)^3<\frac{27}{256}C_{\min}^4
\]
and the decrease of $S(C_*-S)^3$ for $S>C_*/4$ imply $S_i<5$.

The elementary bound
$\log(1+x)\ge x-x^2/[2(1-|x|)]$ for $|x|<1$ and the exact comparison
\[
 \epsilon+\frac9{2p_0(1-3/p_0)}<\frac1{11}
\]
give $(\tau/f_e)^p\le e^{1/11}e^{-z_e}$ on contacts.
For a noncontact, $f_e(0)\ge2\tau$, and the chart derivative bound
$|Df_e(q)[v]|\le4|v|$ implies $f_e>3\tau/2$ on
\eqref{tail:convexset}. Each vertex has three noncontacts. Hence
\begin{equation}\label{tail:incident}
 \tau^{p+1}\sum_{j\ne i}f_{ij}^{-p-1}
 <\frac{5e^{1/11}}{1-3/p_0}+\frac1{1000}<6,
\end{equation}
where $3(2/3)^{p+1}\le3(2/3)^{30}<1/1000$.

For clarity, the derivative estimates used here follow directly from
the chart. On $|q|\le1/2$, $\|Dy_i\|\le(1-t_i)^{-1/2}$ and
$\|D^2y_i\|\le(1-t_i)^{-3/2}$, giving
\[
 |Df_{ij}(q)[v]|\le4|v|,\qquad
 |D^2f_{ij}(q)[v,v]|\le6(|v_i|^2+|v_j|^2).
\]
Integration along a radius and the four contacts at each vertex give
$\|Df_{\mathcal C}(q)-R\|_2\le12|q|$.
Lemma~\ref{geo:spectrum} and the endpoint test therefore imply
\[
 \sum_{e\in\mathcal C}(Df_e(q)[v])^2
 \ge(\sqrt3-12K_f/p_0)^2|v|^2>\frac52|v|^2\quad(v\ne0).
\]
An upper Hessian bound improves the absolute bound on the smaller ball.
Set $r_f=K_f/p_0$ and $h_i=\sqrt{1-t_i}$. Direct differentiation gives
\[
 D^2f_{ij}[v,v]=-4Dy_i[v_i]\cdot Dy_j[v_j]
 +2(P_i\cdot y_j)b_i+2(P_j\cdot y_i)b_j,
\]
where $b_i=|v_i|^2/h_i+(q_i\cdot v_i)^2/h_i^3$.
Since $P_i\cdot P_j\le a$ and
$|y_j-P_j|\le(101/100)r_f$, the coefficient is bounded by
\[
 \frac2{1-r_f^2}+\frac{2(a+(101/100)r_f)}{(1-r_f^2)^{3/2}}<\frac83.
\]
The displacement bound and this comparison are checked at $r_f$;
both remain valid on every smaller ball. Thus, for all 28 pairs,
\begin{equation}\label{tail:pairhessian}
 D^2f_{ij}[v,v]\le\frac83(|v_i|^2+|v_j|^2).
\end{equation}

Differentiate the entire energy:
\[
 D^2E_s(q)[v,v]=p(p+1)\sum_{i<j}f_{ij}^{-p-2}(Df_{ij}[v])^2
                 -p\sum_{i<j}f_{ij}^{-p-1}D^2f_{ij}[v,v].
\]
In the positive sum retain the contact terms and use
\eqref{tail:contactbounds} and the differential bound. In the negative
sum use \eqref{tail:incident} and \eqref{tail:pairhessian} for every
pair, including noncontacts. This proves
\begin{equation}\label{tail:hessian}
 D^2E_s(q)[v,v]\ge p\tau^{-p-2}
       \left[\frac34(p+1)-16\tau\right]|v|^2>0\quad(v\ne0)
\end{equation}
throughout $\mathcal K_p$. The bracket is positive at $p_0$ and
increases with $p$.

\begin{proof}[Proof of Theorem~\ref{thm:main} for $s\ge124$]
Compactness and lower semicontinuity give a global minimizer.
The trial ceiling makes its energy finite, so it is collision-free
and is a local minimum. The closest-pair argument, global packing gap
and radius improvements place an aligned representative in
$\mathcal K_p$. The fully stationary candidate
also has a representative there.
Their joining segment stays in $\mathcal K_p$. If the representatives
were distinct, \eqref{tail:hessian} would make the energy strictly
convex on the segment. Its derivative at the candidate endpoint is
zero, so the other endpoint would have strictly greater energy.
This contradicts global minimality. Restoring orthogonal transformations
and labels proves precisely the stated equality case.
\end{proof}

\section{Verification and arithmetic}\label{sec:arithmetic}
The mathematical proofs above reduce the computation to finite
identities, strict inequalities, and complete interval coverage.
We describe how the programs verify these obligations. Certificate
search supplies numerical data; the proof uses only the resulting
verified identities and inequalities.

\subsection{Uniform verification of a finite cell}
A cell is the exact closed interval $s=c+h\zeta$, $-1\le\zeta\le1$,
with rational $c,h$. The Taylor order is independent of the Hermite
degree. A Taylor model consists of polynomial coefficients and a
remainder interval that encloses the exact represented function over
the entire cell. All arithmetic and root errors enter its remainder.

For the height, let $\widetilde a(s)$ be a proposed approximation.
On a tube around its graph the checker bounds $F_a\ge m_0>0$ and
$|F(s,\widetilde a(s))|\le\eta$, with $\eta/m_0$ smaller than the
tube padding. The intermediate value and mean value theorems then
give the unique exact root and
$|a_s-\widetilde a(s)|\le\eta/m_0$.
The root enclosure must lie in $(0,1/3)$, and its error is propagated
through the nodes, interpolation, kernels and coefficient system.

For the correction \eqref{finite73:correction}, let $P$ and $Z$ be
proposed matrix and vector polynomial models for $T^{-1}$ and
$T^{-1}r$. Bound the complete residuals
\[
 R_T=I-TP,\qquad \rho_r=r-TZ.
\]
If $B\ge|R_T|$ entrywise and $q=\|B\|_\infty<1$, the Neumann series
proves invertibility and
\begin{equation}\label{arith:solve}
 |T^{-1}r-Z|\le |P|(I-B)^{-1}|\rho_r|.
\end{equation}
For a nonnegative vector $r_0\ge|\rho_r|$, the constant vector
$\|r_0\|_\infty/(1-q)$ dominates $(I-B)^{-1}r_0$.
The update $v\mapsto r_0+Bv$ preserves this domination. Thus any
finite number of outward iterations gives a proved bound; convergence
after a prescribed number of steps is unnecessary.
Propagating \eqref{arith:solve} through the correction matrix encloses
the exact corrected positive-definiteness blocks.

For a corrected symmetric block $B_*$ and a square proposed
preconditioner $P_*$, it suffices to prove
\begin{equation}\label{arith:positive}
 \|P_*B_*P_*^T-I\|_2<1.
\end{equation}
The congruence is then positive definite, $P_*$ is nonsingular, and
$B_*$ is positive definite. The square-shape condition is essential.
An outward row-sum norm of an entrywise absolute bound controls the
spectral norm of the symmetric matrix. A sharper coefficient bound
may use a proposed square eigenbasis $Q$: if
$\varepsilon=\|Q^TQ-I\|_\infty<1$, then, for symmetric $M$,
\[
 \|M\|_2\le\frac{\|Q^TMQ\|_\infty}{1-\varepsilon}.
\]
The lower singular-value bound of $Q$ proves this inequality, so an
approximate eigenbasis is sufficient once its residual is enclosed.

The checker uses exact identities to reduce the interval widths.
The free rows of a kernel basis are fixed rows of its constant
change-of-basis matrix; their positive Taylor coefficients and
remainders are exactly zero. A recursively constructed formal inverse
has product coefficients $I,0,\ldots,0$ through its construction order;
its omitted-degree terms and input remainders still have to be bounded.
In contrast, all coefficients of a merely proposed correction inverse
are checked by the complete residual above. Streamed contractions are
partitions of a sum and retain every cross term and remainder.
The higher-precision replacement of a constant residual coefficient
encloses that same exact coefficient, including the height error.
When a stored block is exactly the identity, a free/free contraction
may be performed by exact integer multiplication of the dyadic
change-of-basis entries. The predicate is checked before using this
special case; otherwise the general interval contraction is used.
These operations change the enclosures, not the defining identity.

\subsection{Exact arithmetic}
The packing enumeration uses integers. The geometric and final tail
checkers use rational intervals, exact arithmetic in
$\mathbb Q(\sqrt2)$, integer square-root bounds, and exponential and
logarithmic series with explicit tails. All divisions and square roots
have their sign conditions checked. An interval $LDL^T$ computation
with positive pivot lower bounds certifies positive definiteness of
the exact symmetric matrix it encloses.

The finite checker also proves its bounds by integer arithmetic.
Every finite stored binary64 number denotes its exact dyadic rational.
Scalar interval operations use integer floor and ceiling division;
directed conversion to this storage format is performed on its bits.
The higher-precision layer uses dyadic endpoints with a fixed common
denominator. Reciprocals reject intervals containing zero; nonfinite
inputs and bounds are rejected.

For matrix multiplication, write the input enclosures as
$2^{-p}(M+[-R,R])$ and $2^{-q}(N+[-S,S])$, where $M,N$ are integer
matrices and $R,S$ are nonnegative integer matrices. Their product is
contained in
\[
 2^{-p-q}\bigl(MN+[-H,H]\bigr),\qquad
 H=|M|S+R|N|+RS.
\]
All four products are exact integer products. Outward integer
quantization of the inputs and outputs preserves this enclosure.
Thus floating-point dot-product accuracy is not a proof assumption.

Square roots are bounded by integer square-root inequalities.
For exponentials, repeated halving reduces a nonnegative argument to
$[0,1/8]$; the positive Taylor series and a geometric bound on its
omitted terms give an enclosure, which is propagated by squaring.
Negative arguments use reciprocals. For logarithms, powers of two
reduce the argument to $z\in[1,2]$. The series
\[
 \log z=2\sum_{j\ge0}\frac{u^{2j+1}}{2j+1},
 \qquad u=\frac{z-1}{z+1},
\]
has a geometric tail bound since $0\le u\le1/3$.
All range and sign conditions are checked. Taylor models include all
omitted products and the analytic remainders described above.
Numerical root finding, inverses and eigenbases supply proposals only;
their complete residuals are checked with these exact enclosures.

The raw polynomial tensors are integers. For permitted half-degrees
$6\le d\le10$, the recurrence bounds the coefficient sum of $Q_9$
by $713216$ and the moment numerator magnitude by
$63\cdot6\cdot713216+6=269595654<2^{53}$.
Thus their signed 64-bit storage is exact. The executable checker
still trusts the Python runtime, its integer operations and the FLINT
integer matrix implementation. These programs have not been formally
verified in Lean.

\subsection{Inputs, checks and coverage}
The finite checker and its exact cell registry are in
\path{finite/kernel/} and \path{finite/cells.json}. The witnesses are
in \path{certificates/finite/}. The geometric and tail
certificates are in \path{tail/}.
The program \path{finite/replay.py} checks source and witness hashes,
legal pivots, matrix dimensions, finite entries and square preconditioners.
It verifies the exact cell endpoints, degree, interpolation profile
and Taylor order before aggregating the fresh reports.

The compact witnesses store proposals from which the constant
correction matrix $C$ is reconstructed. Its entries need not satisfy
any asserted reconstruction identity: the actual finite output entries
define an arbitrary dyadic matrix, and the complete uniform residual
test proves its suitability. Small factors replace the dense matrix
in the stored data; optional low-rank terms retain selected search
directions. The lossless encoding and the proposal reconstruction are
specified alongside the checker. A future formal checker may accept
the generated dyadic matrix as an untrusted input in exactly the same way.

For each finite cell the accepted report must establish the root and
node conditions, positive moment pivots, invertible evaluation and
correction systems, strict positivity of every corrected block,
invertibility of the zeroth moment change of basis, and the three
forbidden-triangle exclusions. All endpoints are rational and cells
are closed; the union is recomputed exactly.

All $111$ finite cells passed with the compact witnesses and the exact
arithmetic checker: $98$ cover $[0,74]$ and $13$ cover $[74,124]$.
The source, witness and report bindings are retained in
\path{verification/}; the exact finite union is $[0,124]$.

The five exact tail checkers cover packing enumeration, the signed mean,
contact variance, the quartic packing gap, and the scalar inequalities.
They were independently replayed. The final driver
passes $403$ checks, including all $316$ inner-annulus mass
alternatives and the two outer-annulus endpoint tests. The source and
witness bindings accompany the reports in \path{verification/}.
These computations establish the stated premises using the integer
arithmetic above. Since the finite cells include $124$ and the tail applies
to every $s\ge124$, their union is $[0,\infty)$.

\subsection{Interface for a later formalization}
The proof separates three tasks. First, prove the finite certificate
criterion, the packing and localization reductions, and the convexity
argument as mathematical lemmas. Second, verify the finite arithmetic
obligations for the bound witnesses. Third, verify the rational union
of the exponent cells and the endpoint of the unbounded range.
Certificate search need not be formalized: its output is accepted only
after the checking task proves the required identities and inequalities.
Both arithmetic backends reduce to finite integer computations and
explicit series bounds. A Lean formalization still requires proofs of
the interval and Taylor-model operations and of their mathematical
implications; the present replay is not a Lean certificate.

\begin{thebibliography}{1}
\bibitem{CW}
H.~Cohn and J.~Woo,
\emph{Three-point bounds for energy minimization},
J. Amer. Math. Soc. \textbf{25} (2012), 929--958.
\end{thebibliography}
\end{document}
